\documentclass{situationist-preprint}

\title{Brutalism and Its Critics}
\msid{AES-17-b}
\msfamily{AES-17}
\msdiscipline{Aesthetics}
\mstype{Companion manuscript}
\mskeywords{architecture; Brutalism; public art; captive audience; social housing}

\begin{document}
\maketitle

\begin{abstract}
Brutalist housing estates have been judged very differently by architectural critics and by many of the people who lived in them. Critics have praised their honesty of material, their sculptural force and their social ambition; residents have often found them harsh, oppressive and hard to live with. We argue that the divergence is explained by the captive audience that AES-17-a identifies as the mark of architecture's public character. The critic chooses to attend and judges the building as an artwork; the resident cannot opt out and meets the building as a condition of daily life. The stricter standards of architecture apply to the second encounter, and a building can deserve the critic's praise while failing them.
\end{abstract}

\section{The divergence}

Brutalism, named and defended by \citet{Banham1966}, prized the exposure of raw concrete, the legibility of structure and services, and a monumental scale. Its large housing schemes in Britain and elsewhere, among them the Smithsons' Robin Hood Gardens in London, were admired by critics as serious attempts to give architectural form to collective life. Many residents came to regard such estates as bleak and alienating, and several were demolished while critics campaigned for their preservation.

The divergence is usually explained as a conflict of taste, with the critics accused of elitism or the residents of philistinism. We propose an explanation that requires neither, and that takes seriously the old thought that a building is experienced in use and in passing rather than only in contemplation \citep{Scruton1979}.

\section{Two audiences}

By Definition~2.1 of AES-17-a, a work has a public character when it is permanent, located in shared space, and encountered as a condition of ordinary life by people who have not chosen to attend to it, its captive audience. Principle~P3 of AES-17-a holds that a work's license to impose discomfort or disquiet depends on its audience's having chosen it; where it has not, such responses must be warranted by the work's public purpose.

The critic and the resident stand in different relations to the same building. The critic visits, attends and judges the building as an artwork among others, and for an audience that has chosen to attend, the austerity and severity of Brutalism are licensed in the way the severities of other arts are. The resident is captive. For her the same severity is imposed, daily and without exit.

\begin{proposition}\label{p:split}
A building can possess the aesthetic properties and artistic value for which critics praise it, and also fail the stricter standards that its public character imposes. Where it does, the critics' and residents' judgments are both correct, about different things.
\end{proposition}

\begin{proof}[Argument]
Theorem~3.2 of AES-17-a states that architecture is an art, so a building's artistic value can be assessed as critics assess it. Theorem~3.4 states that architecture is subject to stricter standards (Definition~2.2 of AES-17-a): it must meet the conditions of the other arts and further ones, and is denied excuses they enjoy. One such further condition, by Principle~P3, is that imposed severity be warranted by public purpose. Whether a housing estate meets it depends on whether its severity serves the lives of its residents, a question the residents are best placed to answer. Meeting the first set of conditions does not entail meeting the second.
\end{proof}

\section{Objections}

\begin{objection}
Residents' complaints were often about maintenance, management and social policy rather than architecture.
\end{objection}

\begin{reply}
Many were, and those complaints do not bear on the buildings' standing. The proposition concerns only the part of the residents' judgment that answers to the buildings' perceptible character, which is the part a captive audience is in a position to judge.
\end{reply}

\begin{objection}
Some Brutalist estates later became sought after. This shows that the residents' judgment was a matter of fashion.
\end{objection}

\begin{reply}
A later audience that chooses to live in a building is not captive in the sense of Definition~2.1 of AES-17-a; its judgment confirms the critics' assessment of artistic value and leaves the earlier failure toward a captive audience in place.
\end{reply}

\section{Conclusion}

The quarrel over Brutalist housing is not a quarrel between good and bad taste. It shows architecture's public character at work: one building, judged as art by those who chose it and as a condition of life by those who did not, and held to stricter standards by the second.

\begin{thebibliography}{9}\small
\bibitem[Banham(1966)]{Banham1966} Banham, R. (1966). \emph{The New Brutalism: Ethic or Aesthetic?} Architectural Press.
\bibitem[Scruton(1979)]{Scruton1979} Scruton, R. (1979). \emph{The Aesthetics of Architecture}. Princeton University Press.
\end{thebibliography}

\end{document}
